\clearpage\fancyhead[L]{\small COMET.}\SourceLeaf{387}
\begin{center}\textbf{COMET.}\par\vspace{4pt}\hrule\vspace{4pt}\fontsize{10}{12}\selectfont Göttingische gelehrte Anzeigen. Stück 45. Seite 441. 1815 März 20.\par\vspace{4pt}\hrule\vspace{4pt}\end{center}

\par\indent In einem Schreiben an den Herrn Prof.\ \textsc{Gauss} vom 7.\ März theilte Hr.\ Doctor \textsc{Olbers} die Nachricht mit, dass er am 6.\ März im Perseus einen neuen Cometen entdeckt habe. Den 6.\ März $10^{\mathrm u}\,55^{\mathrm m}$ mittl.\ Zeit war seine gerade Aufsteigung $49^\circ\,7'$, nordliche Abweichung $32^\circ\,7'$. Den 7.\ März $7^{\mathrm u}\,40^{\mathrm m}$ die gerade Aufsteigung $49^\circ\,22'$, nordliche Abweichung $32^\circ\,32'$. Es lässt sich hiernach erwarten, dass dieser Comet noch ziemlich lange sichtbar bleiben werde. Uebrigens ist er sehr klein, hat einen schlecht begrenzten Kern, und einen sehr blassen durchsichtigen Nebel, und ist dem zweiten Cometen von 1811, wie dieser im Februar 1812 erschien, an Licht und Form ähnlich.

\par\indent Gleich nach Eingang dieses Briefes am 13.\ März wurde dieser neue Comet auf der hiesigen Sternwarte aufgesucht. Nur auf wenige Minuten klärte sich der Himmel so weit auf, dass der Comet zwischen 159 und 164 im Perseus (nach \textsc{Bodens} Catalog) bemerkt werden konnte. Allein zu schnell folgten schon wieder Regenwolken, so dass auch nicht einmal eine Schätzung möglich war.
